% Conteudo migrado de thesis/main.tex.
% O titulo do capitulo fica definido em thesis.tex.


\section{Limitações do Sistema Atual}

Embora o sistema de aquisição da equipe seja robusto e consolidado, sua arquitetura
cabeada impõe restrições estruturais que limitam a capacidade de instrumentação
do veículo. Essas limitações se manifestam em três cenários distintos, cada um
com suas próprias implicações práticas.

O primeiro cenário envolve componentes rotativos, como o interior das rodas.
Nesses pontos, a conexão com o chicote elétrico do veículo é tecnicamente
possível por meio de acopladores rotativos, porém a complexidade de instalação,
o custo e a fragilidade desses componentes em ambiente de competição tornam
a solução inviável na prática. Isso impede a leitura de sensores como temperatura
de cubo ou deformação de manga de eixo — informações de alto valor para a dinâmica
veicular.

O segundo cenário diz respeito a regiões onde o cabeamento interfere diretamente
no fenômeno a ser medido ou onde seu custo de instalação não se justifica pela
frequência de uso. A asa traseira é o exemplo mais crítico: trata-se de um componente
aerodinâmico cuja performance é sensível a qualquer perturbação no escoamento
de ar ao redor de sua superfície. Além disso, testes de instrumentação na asa
são pontuais e pouco frequentes, o que torna ainda menos justificável o
investimento em um chicote dedicado — tanto em termos de material quanto de
mão de obra para montagem e desmontagem a cada sessão.

O terceiro cenário extrapola os limites físicos do veículo. Sensores posicionados
fora do carro — como um sensor de passagem para detecção automática de voltas,
ou equipamentos de medição estacionários utilizados em testes controlados —
não têm como se comunicar com o sistema de aquisição embarcado por meios
cabeados. Essa limitação impede a integração de dados externos ao fluxo de aquisição,
obrigando a equipe a recorrer a registros manuais ou sistemas completamente
separados, sem sincronização temporal com os demais canais do carro.
% what are the limitations of our current setup
% where they are more proeminent

